%% Solutions for Chapter 9: Normal Forms and Pumping Lemmas for Context-Free Languages

\chapter{Normal Forms and Pumping Lemmas --- Solutions}

\begin{enumerate}[{9.}1. ]

%%------------------------------------------------------------
\item \textbf{Characterize the parse trees of left-linear grammars.}

In a left-linear grammar, every production has the form $A \to w$ (terminal string) or $A \to Bw$ (nonterminal followed by a terminal string). Consequently, in any parse tree for a derivation in a left-linear grammar:
\begin{itemize}
\item The tree has a ``left-spine'': the leftmost child of any internal node is always a nonterminal, and the remaining children are terminal symbols.
\item All nonterminals appear in the leftmost branch of the parse tree.
\item The tree is maximally left-branching; at every level, the left subtree contains all nonterminal structure while the right siblings are terminals.
\end{itemize}
This is the mirror image of right-linear parse trees (which are maximally right-branching). The depth of the parse tree equals the number of nonterminals in the derivation.

%%------------------------------------------------------------
\item \textbf{Context-free grammars for the following languages.}

\begin{enumerate}
\item $\{\pmb{a}^n\pmb{b}^n\pmb{c}^m\pmb{d}^m \mid n,m \in \NN\}$:
\[
S \to AB, \quad A \to \pmb{a}A\pmb{b} \mid \lambda, \quad B \to \pmb{c}B\pmb{d} \mid \lambda.
\]

\item $\{\pmb{a}^i\pmb{b}^j\pmb{c}^j\pmb{d}^i \mid i,j \in \NN\}$:
\[
S \to \pmb{a}S\pmb{d} \mid T \mid \lambda, \quad T \to \pmb{b}T\pmb{c} \mid \lambda.
\]

\item $\{\pmb{a}^n\pmb{b}^n\pmb{c}^m\pmb{d}^m\} \cup \{\pmb{a}^i\pmb{b}^j\pmb{c}^j\pmb{d}^i\}$:
\[
S \to S_1 \mid S_2
\]
where $S_1$ generates the first language (as in part (a)) and $S_2$ generates the second (as in part (b)).
\end{enumerate}

%%------------------------------------------------------------
\item \textbf{Unambiguous context-free grammars.}

\begin{enumerate}
\item
\begin{enumerate}[i.]
\item For $\{\pmb{a}^n\pmb{b}^n\pmb{c}^m\pmb{d}^m\}$: The grammar $S \to AB$, $A \to \pmb{a}A\pmb{b}\mid\lambda$, $B \to \pmb{c}B\pmb{d}\mid\lambda$ is already unambiguous (every string has a unique parse).

\item For $\{\pmb{a}^i\pmb{b}^j\pmb{c}^j\pmb{d}^i\}$: The grammar $S \to \pmb{a}S\pmb{d}\mid T\mid\lambda$, $T \to \pmb{b}T\pmb{c}\mid\lambda$ is unambiguous.

\item For the union $L_1\cup L_2$, it \emph{is} possible to make the grammar
unambiguous by splitting the union into disjoint pieces:
\[
L_1\cup L_2=L_1\cup\{a^i b^j c^j d^i\mid i<j\}\cup\{a^i b^j c^j d^i\mid i>j\}.
\]
An unambiguous grammar is
\[
S\to S_1\mid S_<\mid S_>,
\]
where $S_1$ is the unambiguous grammar from part (i), and
\[
S_<\to \pmb{a}\pmb{b}S_<\pmb{c}\pmb{d}\mid U,\qquad
U\to \pmb{b}U\pmb{c}\mid \pmb{b}\pmb{c},
\]
\[
S_>\to \pmb{a}S_>\pmb{d}\mid \pmb{a}V\pmb{d},\qquad
V\to \pmb{a}\pmb{b}V\pmb{c}\pmb{d}\mid \lambda.
\]
The $S_<$ branch generates exactly the case $j>i$, the $S_>$ branch generates
exactly the case $i>j$, and the $S_1$ branch generates $L_1$.  These three
subsets are disjoint, so the grammar is unambiguous.
\end{enumerate}

\item The statement is \emph{false}.  The example from part (i)(iii) already
shows that two unambiguous context-free languages can have a union that is
still unambiguous.  On the other hand, Theorem~9.1 shows that inherently
ambiguous context-free languages do exist, and the standard example
\[
\{a^i b^j c^k\mid i=j\}\cup\{a^i b^j c^k\mid j=k\}
\]
shows that the union of two unambiguous context-free languages can fail to be
unambiguous.

\item No.  Since there are unambiguous context-free languages whose union is not
unambiguous, \USIG\ is not closed under union.
\end{enumerate}

%%------------------------------------------------------------
\item \textbf{State and prove the inductive result needed in Theorem 9.2.}

Let $G'=G_{(A_G^{\sigma d})}$ be the grammar built in Theorem~9.2 from the
deterministic automaton $A_G^{\sigma d}$.

\textbf{Inductive statement}: For every state $q$ of $A_G^{\sigma d}$ and every
word $x\in\Sigma^n$, there is at most one derivation in $G'$ from the
nonterminal $q$ that yields $x$.

\textbf{Proof}: By induction on $n$.

\textbf{Base} ($n=0$): The only possible word is $\lambda$.  In the grammar
obtained from a deterministic automaton, the only way to derive $\lambda$ from
$q$ is via the production $q\to\lambda$, which exists exactly when $q$ is a
final state.  Hence there is at most one such derivation.

\textbf{Step}: Assume the statement true for words of length $n$, and let
$x=\pmb{a}y$ with $|y|=n$.  Any derivation of $x$ from $q$ must begin with a
production of the form
\[
q\to \pmb{a}r,
\]
where $r=\delta(q,\pmb{a})$.  Since the machine is deterministic, there is at
most one such state $r$, hence at most one possible first production.  After
that first step, the rest of the derivation must derive $y$ from $r$, and by
the induction hypothesis there is at most one way to do so.  Therefore there is
at most one derivation of $x$ from $q$.

Applying this to the start state shows that every word in $L(G')$ has a unique
derivation tree, so $G'$ is unambiguous. \qed

%%------------------------------------------------------------
\item \textbf{Algorithm to compute $B^u = \{C \mid B \DRightarrow C\}$ for each nonterminal $B$ in $P^u$.}

Consider the unit-production graph: nodes are nonterminals, directed edge $B \to C$ for each unit production $B \to C$ in $P^u$. Then $B^u$ is the set of nonterminals reachable from $B$ in this graph (including $B$ itself).

\textbf{Algorithm}: For each nonterminal $B$, compute $B^u$ by BFS or DFS from $B$ in the unit-production graph:
\begin{enumerate}
\item Initialize $B^u = \{B\}$.
\item While there exists $C \in B^u$ with a unit production $C \to D \in P^u$ and $D \not in B^u$: add $D$ to $B^u$.
\item Repeat until no new nonterminals are added.
\end{enumerate}
This is a standard graph reachability computation and terminates since $|\Omega|$ is finite.

%%------------------------------------------------------------
\item \textbf{Closure of \CSIG\ under concatenation.}

\begin{enumerate}
\item \textbf{Problem with naive construction}: The production $Z \to S_1 \cdot S_2$ is fine as a context-free production, but the issue is with lambda productions. If $\lambda \in L(G_1)$ (i.e., $S_1 \DRightarrow \lambda$), then $Z \to S_1 S_2 \DRightarrow S_2$, so $L(G^*) \supseteq L(G_2)$, not just $L(G_1)\cdot L(G_2)$ minus $\lambda$-contributions. Similarly, if $\lambda \in L(G_2)$, then $Z \DRightarrow S_1$, giving $L(G^*) \supseteq L(G_1)$. The issue: when both $G_1$ and $G_2$ derive $\lambda$, the simple construction may include unintended strings or exclude others due to improper lambda-production treatment.

\item \textbf{Fix}: First transform $G_1$ and $G_2$ so that $\lambda$ is not derivable from their respective start symbols (or handle it explicitly). If $\lambda \in L(G_i)$, add explicit productions:
\[
Z \to S_1 S_2 \mid S_1 \mid S_2 \mid \lambda \quad \text{(appropriate subset based on whether } \lambda \in L(G_i)\text{)}
\]
Specifically: add $Z \to S_2$ if $\lambda \in L(G_1)$; add $Z \to S_1$ if $\lambda \in L(G_2)$; add $Z \to \lambda$ if $\lambda \in L(G_1) \cap L(G_2)$.

\item \textbf{Proof}: $w \in L(G^*_{\text{corrected}})$ iff $Z \DRightarrow w$ iff $w = uv$ with $u \in L(G_1)$ and $v \in L(G_2)$ (considering the $\lambda$ cases via the added productions) iff $w \in L(G_1)\cdot L(G_2)$. \qed
\end{enumerate}

%%------------------------------------------------------------
\item \textbf{$\{\pmb{a}^i\pmb{b}^j\pmb{c}^k \mid i,j,k \in \NN,\ i+j=k\}$ is context free.}

\textbf{Grammar}: $S \to \pmb{a}S\pmb{c} \mid T$, $T \to \pmb{b}T\pmb{c} \mid \lambda$.

\textbf{Proof}: Any derivation produces $\pmb{a}^i \cdot \pmb{b}^j\pmb{c}^j \cdot \pmb{c}^i = \pmb{a}^i\pmb{b}^j\pmb{c}^{i+j}$. Taking $k = i+j$, we get exactly the language. \qed

%%------------------------------------------------------------
\item \textbf{Ambiguous right-linear grammar and its disambiguation.}

\begin{enumerate}
\item $G = \langle\{S,A,B\},\{\pmb{a}\},S,\{S\to A, S\to B, A\to\pmb{aa}A, A\to\lambda, B\to\pmb{aaa}B, B\to\lambda\}\rangle$.
The string $\pmb{aaaaaa} = \pmb{a}^6$ can be derived as $S\Rightarrow A \DRightarrow \pmb{aa}A\DRightarrow \pmb{aaaa}A\DRightarrow \pmb{aaaaaa}$ (via $A$: three $\pmb{aa}$ steps) or as $S \Rightarrow B \DRightarrow \pmb{aaa}B\DRightarrow\pmb{aaaaaa}$ (via $B$: two $\pmb{aaa}$ steps). So $G$ is ambiguous.

\item $L(G) = \{w \mid |w| \equiv 0 \pmod{2}\} \cup \{w \mid |w| \equiv 0 \pmod{3}\} = \{w \mid 2 \mid |w| \text{ or } 3 \mid |w|\}$.
An unambiguous right-linear grammar is obtained from the DFA whose states record
$|w|\bmod 6$.  Let the nonterminals be
$S_0,S_1,S_2,S_3,S_4,S_5$, where $S_i$ corresponds to remainder $i$.  The
start symbol is $S_0$, and the accepting remainders are $0,2,3,4$.  The grammar
is
\[
\begin{aligned}
S_0&\to \lambda\mid \pmb{a}S_1,\\
S_1&\to \pmb{a}S_2,\\
S_2&\to \lambda\mid \pmb{a}S_3,\\
S_3&\to \lambda\mid \pmb{a}S_4,\\
S_4&\to \lambda\mid \pmb{a}S_5,\\
S_5&\to \pmb{a}S_0.
\end{aligned}
\]
This grammar is unambiguous because at each nonterminal there is at most one
production beginning with $\pmb{a}$, so each word determines a unique path
through the six remainders.  A derivation ends with $\lambda$ exactly in the
states whose remainders are divisible by $2$ or $3$, so the grammar generates
precisely $L(G)$.
\end{enumerate}

%%------------------------------------------------------------
\item \textbf{Grammar for regular expressions, modified.}

\begin{enumerate}
\item Grammar that strips outermost parentheses:
\[
R \to \pmb{a}\mid\pmb{b}\mid\pmb{c}\mid\pmb{\epsilon}\mid\pmb{\emptyset}\mid E\pmb{\cdot}E\mid E\pmb{\cup}E\mid R^{\pmb{*}}
\]
where $E \to \pmb{(}R\pmb{)}$ or $E$ is any atomic expression.

\item Grammar allowing extraneous parentheses:
\[
R \to \pmb{(}R\pmb{)} \mid \pmb{a}\mid\pmb{b}\mid\pmb{c}\mid\pmb{\epsilon}\mid\pmb{\emptyset}\mid\pmb{(}R\pmb{\cdot}R\pmb{)}\mid\pmb{(}R\pmb{\cup}R\pmb{)}\mid R^{\pmb{*}}
\]
\end{enumerate}

%%------------------------------------------------------------
\item \textbf{Leftmost and rightmost derivations for $\pmb{((a^*\cdot b)\cup(c\cdot d)^*)}$.}

\begin{enumerate}
\item \textbf{Leftmost derivation} (always expand leftmost nonterminal):
\begin{align*}
R &\Rightarrow \pmb{(}R\pmb{\cup}R\pmb{)} \\
&\Rightarrow \pmb{(}\pmb{(}R\pmb{\cdot}R\pmb{)}\pmb{\cup}R\pmb{)} \\
&\Rightarrow \pmb{((}R^{\pmb{*}}\pmb{\cdot}R\pmb{)}\pmb{\cup}R\pmb{)} \\
&\Rightarrow \pmb{((}\pmb{a}^{\pmb{*}}\pmb{\cdot}R\pmb{)}\pmb{\cup}R\pmb{)} \\
&\Rightarrow \pmb{((a^*\cdot}\pmb{b}\pmb{)}\pmb{\cup}R\pmb{)} \\
&\Rightarrow \pmb{((a^*\cdot b)}\pmb{\cup}\pmb{(}R\pmb{\cdot}R\pmb{)}^{\pmb{*}}\pmb{)} \\
&\Rightarrow \pmb{((a^*\cdot b)\cup(}\pmb{c}\pmb{\cdot}R\pmb{)}^{\pmb{*}}\pmb{)} \\
&\Rightarrow \pmb{((a^*\cdot b)\cup(c\cdot}\pmb{d}\pmb{)}^{\pmb{*}}\pmb{)}
\end{align*}

\item \textbf{Rightmost derivation} (always expand rightmost nonterminal):
\begin{align*}
R &\Rightarrow \pmb{(}R\pmb{\cup}R\pmb{)} \\
&\Rightarrow \pmb{(}R\pmb{\cup}\pmb{(}R\pmb{\cdot}R\pmb{)}^{\pmb{*}}\pmb{)} \\
&\Rightarrow \pmb{(}R\pmb{\cup(}R\pmb{\cdot}\pmb{d}\pmb{)}^{\pmb{*}}\pmb{)} \\
&\Rightarrow \pmb{(}R\pmb{\cup(}\pmb{c}\pmb{\cdot d)}^{\pmb{*}}\pmb{)} \\
&\Rightarrow \pmb{(}\pmb{(}R\pmb{\cdot}R\pmb{)}\pmb{\cup(c\cdot d)}^{\pmb{*}}\pmb{)} \\
&\Rightarrow \pmb{((}R\pmb{\cdot}\pmb{b}\pmb{)\cup(c\cdot d)}^{\pmb{*}}\pmb{)} \\
&\Rightarrow \pmb{((}R^{\pmb{*}}\pmb{\cdot b)\cup(c\cdot d)}^{\pmb{*}}\pmb{)} \\
&\Rightarrow \pmb{((a^*\cdot b)\cup(c\cdot d)}^{\pmb{*}}\pmb{)}
\end{align*}
\end{enumerate}

%%------------------------------------------------------------
\item \textbf{Prove $L(G) = L(G')$ (Theorem 9.5) by induction on derivation steps.}

For every production $A\to\alpha_1\alpha_2\cdots\alpha_n$ of $G$, the grammar
$G'$ contains a replacement set that begins with
$A\to\alpha_1Y_{k1}$ and ends with
$Y_{k,n-2}\to\alpha_{n-1}\alpha_n$ (with terminal symbols replaced by the new
symbols $X_{\pmb{a}}$ when necessary).  Those replacement rules together produce
exactly the same strings as the original rule.

\textbf{Claim}: If $S\DRightarrow_G^{\,m} x$ with $x\in\Sigma^*$, then
$S\DRightarrow_{G'} x$.

\textbf{Proof}: By induction on $m$.

\textbf{Base} ($m=1$): Then $S\to x$ is a single production of $G$.  If it is
already in CNF form, the same production is in $G'$.  Otherwise, the
replacement set for that production derives $x$ in $G'$.

\textbf{Step}: Assume the claim true for derivations of length $m$.  Suppose
\[
S\Rightarrow_G \alpha \DRightarrow_G^{\,m} x.
\]
Replace the first production $S\to\alpha$ by its corresponding replacement set
in $G'$.  This produces the same sentential form $\alpha$, except possibly that
some terminal symbols have been replaced by the new nonterminals
$X_{\pmb{a}}$, each of which immediately derives the corresponding terminal.
After these finitely many extra steps, the remainder of the original derivation
can be copied by the induction hypothesis.  Hence $S\DRightarrow_{G'} x$.

Conversely, every derivation in $G'$ can be compressed to one in $G$: each new
symbol $Y_{ki}$ occurs in only one replacement set, so any maximal block of
steps using $Y_{ki}$ symbols is simply the expanded form of one original
production of $G$.  Replacing each such block by the corresponding production of
$G$ yields a derivation of the same terminal string in $G$.

Therefore $L(G)=L(G')$. \qed

%%------------------------------------------------------------
\item \textbf{For CNF grammar $G$ and $x \in L(G)$, $x \neq \lambda$: number of productions in any derivation of $x$ is $2|x|$.}

\textbf{Proof}: In the chapter's definition of CNF, a nonempty derivation begins
with the start step $Z\to S$.  After that initial step, every production is
either of the form $A\to BC$ or $A\to\pmb{a}$.

If $x$ has length $n$, then exactly $n$ terminal productions are needed, one
for each terminal symbol in $x$.  Before those terminal productions can be
applied, the derivation must create exactly $n$ nonterminals ready to become
those terminals.  Starting from the single nonterminal $S$, each production
$A\to BC$ increases the number of nonterminals by $1$, so exactly $n-1$
productions of the form $A\to BC$ are required.

Thus the total number of productions is
\[
1+(n-1)+n=2n=2|x|,
\]
where the initial $1$ is the production $Z\to S$. \qed

%%------------------------------------------------------------
\item \textbf{Bounds on parse tree depth for CNF grammars.}

\begin{enumerate}
\item \textbf{Lower bound}: For $x \in L(G)$ with $|x| = n > 0$, the parse tree has $n$ leaf nodes (one per terminal). Since each internal node has exactly 2 children (CNF), a complete binary tree with $n$ leaves has depth at least $\lceil \log_2 n \rceil$. Thus the depth of the parse tree is at least $\lceil \log_2 n \rceil$.

\item \textbf{Upper bound}: The parse tree can be maximally skewed (a path): at each level, one child is a leaf (terminal) and the other is the next internal node. This gives depth $n-1$ from the root to the last nonterminal, plus 1 for the final terminal, giving depth $n$. (Or $2n-1$ if we count from a different convention.) Thus the depth is at most $n = |x|$.

\textbf{Summary}: $\lceil \log_2 |x| \rceil \leq \text{depth} \leq |x|$.
\end{enumerate}

%%------------------------------------------------------------
\item \textbf{Convert grammars to Chomsky Normal Form.}

\begin{enumerate}
\item $\langle\{S,B,C\},\{\pmb{a},\pmb{b},\pmb{c}\},S,\{S\to\pmb{a}B,S\to\pmb{ab}C,B\to\pmb{bc},C\to\pmb{c}\}\rangle$:

Introduce terminal nonterminals: $T_a \to \pmb{a}$, $T_b \to \pmb{b}$, $T_c \to \pmb{c}$.
- $S \to \pmb{a}B$ becomes $S \to T_a B$ (already binary CNF).
- $S \to \pmb{ab}C$ becomes $S \to T_a D_1$ where $D_1 \to T_b C$.
- $B \to \pmb{bc}$ becomes $B \to T_b T_c$.
- $C \to \pmb{c}$ becomes $C \to T_c$ --- but $C \to T_c$ is a unit production! Replace: wherever $C$ appears as a target, substitute. Since $C \to \pmb{c}$ and $C$ does not appear in the right side of $B$ or $S$ (only in $S \to \pmb{ab}C$), we can just keep $C \to T_c$ and note it will be eliminated: replace $S \to T_a D_1$ (where $D_1 \to T_b C$) by substituting $C$: $D_1 \to T_b T_c$.

Final CNF:
\[
S \to T_a B \mid T_a D_1, \quad D_1 \to T_b T_c, \quad B \to T_b T_c, \quad T_a \to \pmb{a}, \quad T_b \to \pmb{b}, \quad T_c \to \pmb{c}.
\]

\item $\langle\{S,A,B\},\{\pmb{a},\pmb{b},\pmb{c}\},S,\{S\to\pmb{c}BA,S\to B,A\to\pmb{c}B,A\to A\pmb{bb}S,B\to\pmb{aaa}\}\rangle$:

Step 1: Remove unit production $S \to B$: add all $B$-rules to $S$: $S \to \pmb{aaa}$.
Step 2: Introduce terminal symbols: $T_a \to \pmb{a}$, $T_b \to \pmb{b}$, $T_c \to \pmb{c}$.
Step 3: Break long right-hand sides:
- $S \to \pmb{c}BA$: $S \to T_c E_1$, $E_1 \to BA$.
- $S \to \pmb{aaa}$: $S \to T_a E_2$, $E_2 \to T_a T_a$.
- $A \to \pmb{c}B$: $A \to T_c B$.
- $A \to A\pmb{bb}S$: $A \to A E_3$, $E_3 \to T_b E_4$, $E_4 \to T_b S$.
- $B \to \pmb{aaa}$: $B \to T_a E_2$ (reuse $E_2$).

Final CNF:
\begin{align*}
S &\to T_c E_1 \mid T_a E_2 \\
E_1 &\to BA,\quad E_2 \to T_a T_a \\
A &\to T_c B \mid A E_3 \\
E_3 &\to T_b E_4,\quad E_4 \to T_b S \\
B &\to T_a E_2 \\
T_a &\to \pmb{a},\quad T_b \to \pmb{b},\quad T_c \to \pmb{c}
\end{align*}

\item $\langle\{R\},\{\pmb{a},\pmb{b},\pmb{c},\pmb{(},\pmb{)},\pmb{\epsilon},\pmb{\emptyset},\pmb{\cup},\pmb{\cdot},\pmb{^*}\},R,\{R\to\pmb{a}|\pmb{b}|\pmb{c}|\pmb{\epsilon}|\pmb{\emptyset}|(R\pmb{\cdot}R)|(R\pmb{\cup}R)|R^{\pmb{*}}\}\rangle$:

Introduce terminal nonterminals for each terminal symbol: $T_a,T_b,T_c,T_\epsilon,T_\emptyset,T_{(},T_{)},T_\cup,T_\cdot,T_*$.

$R \to \pmb{a}|\pmb{b}|\pmb{c}|\pmb{\epsilon}|\pmb{\emptyset}$: Each of these is a single-terminal production $R \to a$ (with $|a|=1$), which is already in CNF form.

$R \to (R\pmb{\cdot}R)$: Break as $R \to T_{ (} D_1$, $D_1 \to RD_2$, $D_2 \to T_\cdot D_3$, $D_3 \to RT_{)}$.
$R \to (R\pmb{\cup}R)$: Similarly $R \to T_{(}D_4$, $D_4 \to RD_5$, $D_5 \to T_\cup D_6$, $D_6 \to RT_{)}$.
$R \to R^{\pmb{*}}$: $R \to RD_7$, $D_7 \to T_*$---but $D_7 \to T_*$ is a unit production; substitute: $R \to RT_*$.

Final CNF:
\begin{gather*}
R \to \pmb{a}|\pmb{b}|\pmb{c}|\pmb{\epsilon}|\pmb{\emptyset} \mid T_{(}D_1 \mid T_{(}D_4 \mid RT_* \\
D_1 \to RD_2,\quad D_2 \to T_\cdot D_3,\quad D_3 \to RT_{)} \\
D_4 \to RD_5,\quad D_5 \to T_\cup D_6,\quad D_6 \to RT_{)} \\
T_{(} \to \pmb{(},\quad T_{)} \to \pmb{)},\quad T_\cdot \to \pmb{\cdot},\quad T_\cup \to \pmb{\cup},\quad T_* \to \pmb{^*}
\end{gather*}

\item $\langle\{T\},\{\pmb{a},\pmb{b},\pmb{c},\pmb{d},\pmb{-},\pmb{+}\},T,\{T\to\pmb{a}|\pmb{b}|\pmb{c}|\pmb{d}|T\pmb{-}T|T\pmb{+}T\}\rangle$:

$T \to \pmb{a}|\pmb{b}|\pmb{c}|\pmb{d}$: Already in CNF (terminal productions).
$T \to T\pmb{-}T$: $T \to TD_-$ where $D_- \to T_{-}T$ and $T_- \to \pmb{-}$.
$T \to T\pmb{+}T$: $T \to TD_+$ where $D_+ \to T_+T$ and $T_+ \to \pmb{+}$.

Final CNF:
\[
T \to \pmb{a}|\pmb{b}|\pmb{c}|\pmb{d}\mid TD_-\mid TD_+,\quad D_- \to T_-T,\quad D_+ \to T_+T,\quad T_- \to \pmb{-},\quad T_+ \to \pmb{+}.
\]
\end{enumerate}

%%------------------------------------------------------------
\item \textbf{Convert grammars to Greibach Normal Form.}

\begin{enumerate}
\item $\langle\{S_1,S_2\},\{\pmb{a},\pmb{b},\pmb{c},\pmb{d},\pmb{e}\},S_1,\{S_1\to S_2S_1\pmb{e},S_1\to S_2\pmb{b},S_2\to S_1S_2,S_2\to\pmb{c}\}\rangle$:

One equivalent GNF grammar is:
\[
\begin{aligned}
S_1&\to \pmb{c}S_1\pmb{e}\mid \pmb{c}\pmb{b}\mid \pmb{c}S_1\pmb{e}R\mid \pmb{c}\pmb{b}R,\\
S_2&\to \pmb{c}\mid \pmb{c}S_1\pmb{e}S_2\mid \pmb{c}\pmb{b}S_2\mid \pmb{c}S_1\pmb{e}RS_2\mid \pmb{c}\pmb{b}RS_2,\\
R&\to \pmb{c}S_1\pmb{e}\mid \pmb{c}\pmb{b}\mid \pmb{c}S_1\pmb{e}R\mid \pmb{c}\pmb{b}R\\
&\qquad\mid \pmb{c}S_1\pmb{e}S_2S_1\pmb{e}\mid \pmb{c}\pmb{b}S_2S_1\pmb{e}
\mid \pmb{c}S_1\pmb{e}RS_2S_1\pmb{e}\mid \pmb{c}\pmb{b}RS_2S_1\pmb{e}\\
&\qquad\mid \pmb{c}S_1\pmb{e}S_2\pmb{b}\mid \pmb{c}\pmb{b}S_2\pmb{b}
\mid \pmb{c}S_1\pmb{e}RS_2\pmb{b}\mid \pmb{c}\pmb{b}RS_2\pmb{b}.
\end{aligned}
\]
Every production now begins with the terminal $\pmb{c}$, so this grammar is in
GNF.

\item $\langle\{S_1,S_2,S_3\},\ldots,\{S_1\to S_3S_1, S_1\to S_2\pmb{a},S_2\to\pmb{be},S_3\to S_2\pmb{c}\}\rangle$:

$S_2 \to \pmb{be}$ (already in GNF).
$S_3 \to S_2\pmb{c} \to \pmb{be}\pmb{c}$ (substitute): $S_3 \to \pmb{bec}$.
$S_1 \to S_3 S_1 \to \pmb{bec}S_1$ and $S_1 \to S_2\pmb{a} \to \pmb{be}\pmb{a} = \pmb{bea}$.
GNF: $S_1 \to \pmb{bec}S_1 \mid \pmb{bea}$, $S_2 \to \pmb{be}$, $S_3 \to \pmb{bec}$.

\item $\langle\{S_1,S_2,S_3\},\ldots,\{S_1\to S_1S_2\pmb{c},S_1\to\pmb{d}S_3,S_2\to S_1S_1,S_2\to\pmb{a},S_3\to S_3\pmb{e}\}\rangle$:

As printed, this grammar generates no terminal strings at all.  The nonterminal
$S_3$ has only the production $S_3\to S_3\pmb{e}$, so $S_3$ is unproductive.
Hence the production $S_1\to \pmb{d}S_3$ cannot lead to a terminal word.  The
other production for $S_1$ is $S_1\to S_1S_2\pmb{c}$, which still contains
$S_1$ on the left, so it cannot terminate either.  Therefore
\[
L(G)=\emptyset.
\]
An equivalent grammar in pure Greibach normal form is simply
\[
\langle\{S_1\},\{\pmb{a},\pmb{b},\pmb{c},\pmb{d},\pmb{e}\},S_1,\emptyset\rangle.
\]
\end{enumerate}

%%------------------------------------------------------------
\item \textbf{Prove there exists a CFG $G''$ equivalent to $G'$ (with added $A\to\lambda$ rules).}

$G'$ is obtained from $G$ by adding rules $A \to \lambda$ for various
nonterminals $A$. We want to show that there is an equivalent grammar $G''$ in
which all of those lambda productions have been removed, except possibly the
special start rule $Z\to\lambda$.

\textbf{Proof}: Apply the standard lambda-elimination construction to $G'$.
First identify all \emph{nullable} nonterminals, that is, all $A$ for which
$A\DRightarrow^* \lambda$ in $G'$.  For each production in which nullable
nonterminals occur on the right, add all productions obtained by deleting any
chosen subset of those nullable occurrences, except that we do not create an
empty right-hand side unless it is the special start production $Z\to\lambda$.
After all such replacement productions have been added, delete the old
productions $A\to\lambda$.

Every derivation in $G'$ can now be simulated in the new grammar $G''$: whenever
$G'$ uses a nullable nonterminal and later erases it by $A\to\lambda$, the
grammar $G''$ simply chooses the corresponding shortened production from the
start.  Conversely, every shortened production of $G''$ represents a derivation
in $G'$ in which the omitted nullable nonterminals are first produced and then
erased.  Thus $G''$ generates exactly the same language as $G'$, except that the
only remaining lambda production, if needed, is the allowable special rule
$Z\to\lambda$.

Therefore an equivalent grammar $G''$ always exists. \qed

%%------------------------------------------------------------
\item \textbf{Generalization of Lemma 9.1.}

\textbf{Lemma}: If $G = \langle\Omega,\Sigma,S,P\rangle$, $X \to \gamma B\alpha \in P$, all $B$-rules are $\{B \to \beta_i\}$, and $G'$ replaces $X \to \gamma B\alpha$ with $\{X_B \to \gamma\beta_i\alpha\}$, then $L(G) = L(G')$.

\textbf{Proof}: By induction on derivation length.

($L(G) \subseteq L(G')$): Any derivation in $G$ that uses $X \to \gamma B\alpha$ followed later by $B \to \beta_i$ can be replaced in $G'$ by the single production $X_B \to \gamma\beta_i\alpha$. All other derivations are unchanged.

($L(G') \subseteq L(G)$): Any production $X_B \to \gamma\beta_i\alpha$ in $G'$ can be simulated in $G$ by $X \to \gamma B\alpha$ followed by $B \to \beta_i$. \qed

%%------------------------------------------------------------
\item \textbf{$P = \{d^p \mid p \text{ prime}\}$ is not context free.}

\begin{enumerate}
\item \textbf{Pumping theorem proof}: Assume $P$ is context free with pumping constant $n$. Take $s = \pmb{d}^p$ where $p > n$ is prime. By the pumping lemma, $s = uvwxy$ with $|vwx| \leq n$, $|vx| \geq 1$, and $uv^i wx^i y \in P$ for all $i \geq 0$.

Let $|vx| = m$ where $1 \leq m \leq n$. Then $uv^i wx^i y = \pmb{d}^{p + (i-1)m}$. For $i=0$: $\pmb{d}^{p-m}$. For $i=2$: $\pmb{d}^{p+m}$.

Choose $i = p+1$: the resulting string has length $p + pm = p(1+m)$. Since $p \geq 2$ and $m \geq 1$, $p(1+m)$ is composite (product of two integers each $\geq 2$). So $\pmb{d}^{p(1+m)} \not in P$, contradiction. \qed

\item \textbf{Using nonregularity}: The language $P$ is a language over the
single-letter alphabet $\{\pmb{d}\}$.  If $P$ were context free, then
Theorem~9.15 would imply that $P$ is regular, because every context-free
language over a one-letter alphabet is regular.  But $P$ is known to be
nonregular.  This contradiction shows that $P$ is not context free. \qed
\end{enumerate}

%%------------------------------------------------------------
\item \textbf{$\Gamma = \{w\pmb{2}w \mid w \in \{\pmb{0},\pmb{1}\}^*\}$ is not context free.}

Assume $\Gamma$ is context free with pumping constant $n$. Choose $s = \pmb{0}^n\pmb{1}^n\pmb{2}\pmb{0}^n\pmb{1}^n \in \Gamma$ (take $w = \pmb{0}^n\pmb{1}^n$, $|s| = 4n+1$).

By the pumping lemma: $s = uvwxy$, $|vwx| \leq n$, $|vx| \geq 1$, all pumped strings in $\Gamma$.

Since $|vwx|\leq n$, the substring $vwx$ can meet at most one boundary of the
five-block word
\[
\pmb{0}^n\mid \pmb{1}^n\mid \pmb{2}\mid \pmb{0}^n\mid \pmb{1}^n.
\]
We consider the possibilities.

\begin{itemize}
\item If $vwx$ does \emph{not} contain the separator $\pmb{2}$, then it lies
entirely in the left half $\pmb{0}^n\pmb{1}^n$ or entirely in the right half
$\pmb{0}^n\pmb{1}^n$, possibly crossing the internal boundary between
$\pmb{0}$'s and $\pmb{1}$'s.  Pumping changes one side of the separator but
leaves the other side unchanged.  Therefore the resulting word cannot be of the
form $w'\pmb{2}w'$.

\item If $vwx$ does contain the separator $\pmb{2}$, then there is only one
occurrence of $\pmb{2}$ in $s$.  If $\pmb{2}$ appears in $vx$, then pumping with
$i=0$ or $i=2$ changes the number of $\pmb{2}$'s, so the pumped word is not in
$\Gamma$.

\item The only remaining possibility is that $\pmb{2}$ lies in $w$ but not in
$vx$.  Then pumping changes symbols on one side of $\pmb{2}$ while leaving the
other side and the separator itself unchanged.  Again the two halves around the
separator are no longer identical, so the pumped word is not in $\Gamma$.
\end{itemize}

In every case some pumped word lies outside $\Gamma$, contradicting the pumping
lemma.  Therefore $\Gamma$ is not context free. \qed

%%------------------------------------------------------------
\item \textbf{$\Psi = \{ww \mid w \in \{\pmb{0},\pmb{1}\}^*\}$ is not context free.}

Assume $\Psi$ is context free with pumping constant $n$. Choose $s = \pmb{0}^n\pmb{1}^n\pmb{0}^n\pmb{1}^n \in \Psi$ (take $w = \pmb{0}^n\pmb{1}^n$). $|s| = 4n$.

By the pumping lemma: $s = uvwxy$, $|vwx| \leq n$, $|vx| \geq 1$, all pumped strings in $\Psi$.

Since $|vwx| \leq n$, $vwx$ is contained within a span of $n$ consecutive characters. The string $s$ has the pattern $0^n1^n0^n1^n$. The middle two blocks meet at position $2n$; the window $vwx$ cannot span from the first $0^n$-block to the second $0^n$-block.

In all cases, pumping (or de-pumping with $i=0$) changes the length or composition of one half of $s$ without matching change in the other half, so the result cannot have the form $w'w'$. Contradiction. \qed

%%------------------------------------------------------------
\item \textbf{$\Xi = \{\pmb{b}^{2^j} \mid j \in \NN\}$ is not context free.}

Assume $\Xi$ is context free with pumping constant $n$. Choose $j$ such that $2^j > n$, and let $s = \pmb{b}^{2^j}$.

By the pumping lemma: $s = uvwxy$, $|vwx| \leq n$, $|vx| = m \geq 1$. Then $|uv^i wx^i y| = 2^j + (i-1)m$ for each $i \geq 0$.

For $i=2$: length $= 2^j + m$. We need $2^j + m$ to be a power of 2. But $2^j < 2^j + m \leq 2^j + n < 2^{j+1}$ (since $m \leq n < 2^j$). So $2^j + m$ is strictly between $2^j$ and $2^{j+1}$, which means it is not a power of 2. Contradiction. \qed

%%------------------------------------------------------------
\item \textbf{$\Phi = \{\pmb{a}^{j^2} \mid j \in \NN\}$ is not context free; application to $\Phi'$.}

\begin{enumerate}
\item \textbf{$\Phi$ is not context free}: Assume $\Phi$ is context free with pumping constant $n$. Let $j > n$; take $s = \pmb{a}^{j^2}$.

By the pumping lemma: $s = uvwxy$, $|vwx| \leq n$, $|vx| = m \geq 1$. Then pumped strings have length $j^2 + (i-1)m$. For $i=2$: length $j^2 + m$. We need $j^2 + m$ to be a perfect square. Choose $j > n/2$ so that $n < 2j$; then $j^2 < j^2 + m \leq j^2 + n < j^2 + 2j + 1 = (j+1)^2$. So $j^2 + m$ is strictly between two consecutive perfect squares, hence not a perfect square. Contradiction. \qed

\item \textbf{$\Phi'$ not context free via homomorphism}: Define $h: \{\pmb{b},\pmb{c},\pmb{d}\}^* \to \{\pmb{a}\}^*$ by $h(\pmb{c}) = \pmb{a}$ and $h(\pmb{b}) = h(\pmb{d}) = \lambda$. If $\Phi'$ were a CFL, then $h(\Phi') = \Phi$ would also be a CFL (CFLs are closed under homomorphism), contradicting part~(a). Hence $\Phi' \not in \CSIG$.

\item \textbf{Direct Ogden's lemma proof}: Assume $\Phi'$ is context free, and
let $p$ be the constant from Ogden's lemma.  Choose the word
\[
s=\pmb{b}\pmb{c}^{p^2}\pmb{d}^p \in \Phi',
\]
and mark all $p^2$ occurrences of $\pmb{c}$.  Ogden's lemma gives a
decomposition $s=uvwxy$ such that $vx$ contains at least one marked position,
and $vwx$ contains at most $p$ marked positions.

Let $\alpha=|vx|_{\pmb{c}}$ and $\beta=|vx|_{\pmb{d}}$.  Because $vx$ contains a
marked position, $\alpha\geq 1$.  Also $\alpha\leq p$ and $\beta\leq p$.  Pump
once upward, to $uv^2wx^2y$.  Its number of $\pmb{c}$'s is $p^2+\alpha$, while
its number of $\pmb{d}$'s is $p+\beta$.

If $\beta=0$, then the number of $\pmb{d}$'s stays equal to $p$ while the
number of $\pmb{c}$'s changes, so the relation
$|x|_{\pmb{c}}=(|x|_{\pmb{d}})^2$ fails.

If $\beta\geq 1$, then
\[
(p+\beta)^2-p^2=2p\beta+\beta^2 > p \geq \alpha,
\]
so
\[
p^2+\alpha < (p+\beta)^2.
\]
Again the relation $|x|_{\pmb{c}}=(|x|_{\pmb{d}})^2$ fails.  In either case,
$uv^2wx^2y\not in\Phi'$, contradicting Ogden's lemma.  Therefore
$\Phi'$ is not context free.

\item \textbf{Pumping theorem}: \emph{Yes}.  The ordinary pumping theorem also
works.  Assume $\Phi'$ is context free, and let $n$ be its pumping constant.
Take
\[
s=\pmb{b}\pmb{c}^{n^2}\pmb{d}^n \in \Phi'.
\]
Write $s=uvwxy$ with $|vwx|\leq n$ and $|vx|\geq 1$.

If $vx$ contains the initial $\pmb{b}$, then pumping down ($i=0$) removes that
only $\pmb{b}$, so $uv^0wx^0y\not in\Phi'$.

Otherwise, $vx$ lies entirely in the $\pmb{c}^*\pmb{d}^*$ portion of the word.
Let $\alpha=|vx|_{\pmb{c}}$ and $\beta=|vx|_{\pmb{d}}$.  Then
$\alpha+\beta\geq 1$ and $\alpha+\beta\leq n$.

If $\alpha=0$, pumping up changes the number of $\pmb{d}$'s while leaving the
number of $\pmb{c}$'s fixed, so the quadratic relation fails.

If $\beta=0$, pumping up changes the number of $\pmb{c}$'s while leaving the
number of $\pmb{d}$'s fixed, so the quadratic relation fails.

If both $\alpha$ and $\beta$ are positive, then pumping up gives
$n^2+\alpha$ occurrences of $\pmb{c}$ and $n+\beta$ occurrences of
$\pmb{d}$.  But
\[
(n+\beta)^2-n^2=2n\beta+\beta^2 > n \geq \alpha,
\]
so $n^2+\alpha \neq (n+\beta)^2$.  Thus the pumped word is not in $\Phi'$.

In every case some pumped word lies outside $\Phi'$, contradicting the pumping
theorem.  So it is possible to prove directly from the pumping theorem that
$\Phi'$ is not context free.
\end{enumerate}

%%------------------------------------------------------------
\item \textbf{$L = \{y \in \{\pmb{0},\pmb{1}\}^* \mid |y|_{\pmb{0}} = |y|_{\pmb{1}}\}$ is context free.}

\textbf{Grammar}: $S \to \pmb{0}S\pmb{1} \mid \pmb{1}S\pmb{0} \mid SS \mid \lambda$.

\textbf{Proof}: By induction, any string with equal numbers of $\pmb{0}$s and $\pmb{1}$s is derivable, and all derivations yield strings with equal counts. So $L(G) = L$. \qed

%%------------------------------------------------------------
\item \textbf{Formal recursive definition of the path in Theorem 9.9 (pumping lemma).}

\begin{enumerate}
\item Let the parse tree have root $r$ corresponding to start symbol $S$. Define a \emph{leftmost longest path} recursively:
  - \textbf{Boundary}: If $r$ is a leaf, the path is $\{r\}$.
  - \textbf{Rule}: From a node $v$ with children $c_1, c_2, \ldots, c_k$, choose $c_i$ to be the child whose subtree has the greatest number of leaves. (Break ties by choosing the leftmost.) Append $c_i$ to the path, and recurse from $c_i$.

\item The depth of this path is at least $\log_2$ (number of leaves) $= \log_2 |x|$ (by the binary branching property of CNF). If $|x| \geq b^{h+1}$ where $b$ is the maximum branching factor and $h = |\Omega|$, then the path has length $> h$, so some nonterminal repeats on the path. Let $A$ be the repeated nonterminal at positions $i < j$ on the path. Then the subtree at position $j$ generates $w$, and the subtree at position $i$ generates $vwx$ (with $v,x$ possibly empty but $|vx|\geq 1$ by the choice that $vwx \neq w$). Pumping follows from substituting the subtree at $j$ for the subtree at $i$. \qed
\end{enumerate}

%%------------------------------------------------------------
\item \textbf{$\CSIG$ is closed under $\cup$ via direct grammar construction.}

Let $G_1 = \langle\Omega_1,\Sigma,S_1,P_1\rangle$ and $G_2 = \langle\Omega_2,\Sigma,S_2,P_2\rangle$ with $\Omega_1\cap\Omega_2=\emptyset$. Choose $Z \not in \Omega_1\cup\Omega_2$ and define:
\[
G^\cup = \langle\Omega_1\cup\Omega_2\cup\{Z\},\Sigma,Z,P_1\cup P_2\cup\{Z\to S_1,Z\to S_2\}\rangle.
\]
Then $w \in L(G^\cup)$ iff $Z \DRightarrow w$ iff $Z \Rightarrow S_i \DRightarrow w$ for $i=1$ or $i=2$ iff $w \in L(G_1)\cup L(G_2)$. \qed

%%------------------------------------------------------------
\item \textbf{Closure properties of $\XSIG$ (languages that are NOT context free).}

\begin{enumerate}
\item \textbf{Union}: Not closed.  Let
\[
A=\{\pmb{a}^{2^j}\mid j\in\NN\}.
\]
Exercise~9.21 shows that $A$ is not context free.  Let
$B=\pmb{a}^*-A$.  If $B$ were context free, then by Theorem~9.15 it would be
regular, and therefore $A=\pmb{a}^*-B$ would also be regular, contrary to
Exercise~9.21.  Hence $B\in\XSIG$ as well.  But
\[
A\cup B=\pmb{a}^*,
\]
which is regular, hence context free.  Therefore \XSIG\ is not closed under
union.

\item \textbf{Complement}: Not closed.  Lemma~9.4 gives a context-free language
$L$ whose complement is not context free.  Let $X=\overline{L}$.  Then
$X\in\XSIG$, but $\overline{X}=L\in\CSIG$.  So \XSIG\ is not closed under
complement.

\item \textbf{Intersection}: Not closed.  Using the same unary languages $A$ and
$B=\pmb{a}^*-A$ from part (a), both $A$ and $B$ belong to \XSIG, but
\[
A\cap B=\emptyset,
\]
which is context free.  Hence \XSIG\ is not closed under intersection.

\item \textbf{Kleene closure}: Not closed.  For the same language
$A=\{\pmb{a}^{2^j}\mid j\in\NN\}$, we have $A\in\XSIG$.  Since
$\pmb{a}=\pmb{a}^{2^0}\in A$, every unary word belongs to $A^*$, and therefore
\[
A^*=\pmb{a}^*.
\]
Thus $A^*$ is regular, so \XSIG\ is not closed under Kleene closure.

\item \textbf{Concatenation}: Not closed.  Let
\[
A' = A\cup\{\lambda\},\qquad B' = B\cup\{\lambda\},
\]
where $A$ and $B$ are as in part (a).  If $A'$ were context free, then by
Theorem~9.15 it would be regular; removing the single word $\lambda$ would show
that $A$ is regular, contradiction.  So $A'\in\XSIG$, and similarly
$B'\in\XSIG$.

Now every word $\pmb{a}^n$ belongs to $A'B'$: if $\pmb{a}^n\in A$, then
$\pmb{a}^n=\pmb{a}^n\lambda$ with $\pmb{a}^n\in A'$ and $\lambda\in B'$; if
$\pmb{a}^n\not in A$, then $\pmb{a}^n\in B$ and
$\pmb{a}^n=\lambda\pmb{a}^n$ with $\lambda\in A'$ and $\pmb{a}^n\in B'$.
Hence $A'B'=\pmb{a}^*$, which is context free.  Therefore \XSIG\ is not closed
under concatenation.
\end{enumerate}

%%------------------------------------------------------------
\item \textbf{$\CSIG$ is closed under reversal.}

Let $G = \langle\Omega,\Sigma,S,P\rangle$ generate $L$. Define $G^r = \langle\Omega,\Sigma,S,P^r\rangle$ where $P^r = \{A \to \alpha^r \mid A \to \alpha \in P\}$ (reverse the right-hand side of each production).

\textbf{Claim}: $L(G^r) = L^r$.

\textbf{Proof}: We prove a stronger statement: for every nonterminal
$A\in\Omega$ and every terminal string $x\in\Sigma^*$,
\[
A\DRightarrow_G^* x \quad\Longleftrightarrow\quad A\DRightarrow_{G^r}^* x^r.
\]

We prove the forward implication by induction on the derivation length of
$A\DRightarrow_G^* x$.

\textbf{Base}: If the derivation has one step, then $A\to x$ is a production of
$G$.  By definition, $A\to x^r$ is a production of $G^r$, so
$A\DRightarrow_{G^r} x^r$.

\textbf{Step}: Suppose the first production is
\[
A\to X_1X_2\cdots X_k
\]
in $G$, where each $X_i$ is either a terminal or a nonterminal, and suppose the
rest of the derivation yields
\[
X_1X_2\cdots X_k \DRightarrow_G^* x_1x_2\cdots x_k = x,
\]
where each $x_i$ is the terminal string ultimately produced from $X_i$
(possibly a single terminal if $X_i\in\Sigma$).  By the induction hypothesis,
each nonterminal $X_i$ derives $x_i^r$ in $G^r$, while terminals already equal
their own one-letter reversals.  Since $G^r$ contains the production
\[
A\to X_kX_{k-1}\cdots X_1,
\]
we obtain
\[
A \Rightarrow_{G^r} X_kX_{k-1}\cdots X_1 \DRightarrow_{G^r}^*
x_k^r x_{k-1}^r \cdots x_1^r = (x_1x_2\cdots x_k)^r = x^r.
\]

The reverse implication follows by symmetry, since $(G^r)^r = G$.  Taking
$A=S$, we conclude that
\[
x\in L(G) \quad\Longleftrightarrow\quad x^r\in L(G^r),
\]
so $L(G^r)=L^r$. \qed

%%------------------------------------------------------------
\item \textbf{$\CSIG$ is closed under operator $T$ ($T(L) = \{w^r w \mid w \in L\}$).}

\textbf{Claim}: $T$ does not preserve $\CSIG$ in general; i.e., there exists $L\in\CSIG$ such that $T(L)\not in\CSIG$.

\textbf{Counterexample}: Let
\[
L=\{\pmb{a}^n\pmb{b}\pmb{a}^n\mid n\geq 0\}.
\]
This language is context free; for example, it is generated by
$S\to\pmb{a}S\pmb{a}\mid\pmb{b}$.  Since every word of $L$ is a palindrome, we
have
\[
T(L)=\{\pmb{a}^n\pmb{b}\pmb{a}^{2n}\pmb{b}\pmb{a}^n\mid n\geq 0\}.
\]
We claim $T(L)\not in\CSIG$.  Suppose instead that $T(L)$ is context free, and
let $N$ be its pumping constant.  Choose
\[
s=\pmb{a}^N\pmb{b}\pmb{a}^{2N}\pmb{b}\pmb{a}^N \in T(L).
\]
Write $s=uvwxy$ with $|vwx|\leq N$ and $|vx|\geq 1$.

Because $|vwx|\leq N$, the substring $vwx$ can contain at most one of the two
$\pmb{b}$'s.  Thus it lies entirely within one block, or else it crosses just
one boundary between adjacent blocks.

\begin{itemize}
\item If $vwx$ lies entirely within the leftmost $\pmb{a}^N$ block, then pumping
changes only that block.  The rightmost $\pmb{a}$-block remains of length $N$,
so the pumped word cannot have the form
$\pmb{a}^n\pmb{b}\pmb{a}^{2n}\pmb{b}\pmb{a}^n$.

\item If $vwx$ lies entirely within the middle $\pmb{a}^{2N}$ block, then
pumping changes only the middle block.  The outer blocks still both have length
$N$, so the middle block should still have length $2N$, but it does not.

\item If $vwx$ lies entirely within the rightmost $\pmb{a}^N$ block, the
argument is symmetric to the first case.

\item If $vwx$ contains one of the $\pmb{b}$'s, then pumping with $i=0$ or
$i=2$ either changes the number of $\pmb{b}$'s or changes one adjacent
$\pmb{a}$-block while leaving the other outer $\pmb{a}$-block unchanged.  In
either event, the pumped word is not of the form
$\pmb{a}^n\pmb{b}\pmb{a}^{2n}\pmb{b}\pmb{a}^n$.
\end{itemize}

Every possible placement of $vwx$ leads to a contradiction.  Therefore
$T(L)$ is not context free, and \CSIG\ is not closed under $T$. \qed

%%------------------------------------------------------------
\item \textbf{$\mathfrak{C}_{\{\pmb{a},\pmb{b}\}}$ is not closed under relative complement.}

\textbf{Disproof}: Let $L_1 = \{x \in \{\pmb{a},\pmb{b}\}^* \mid |x|_{\pmb{a}} = |x|_{\pmb{b}}\}$ (CFL) and $L_2 = \{\pmb{a}\}^*\{\pmb{b}\}^*$ (regular, hence CFL). Then:
\[
L_1 - L_2 = \{x \mid |x|_{\pmb{a}} = |x|_{\pmb{b}}\} - \{\pmb{a}^n\pmb{b}^n\} = \{x \mid |x|_{\pmb{a}}=|x|_{\pmb{b}}, x \not in \pmb{a}^*\pmb{b}^*\}.
\]
This language (balanced strings that are not in $\pmb{a}^*\pmb{b}^*$) is
context free (construct a grammar). So this does not give a counterexample.

A better approach: use the fact that if \CSIG\ were closed under relative
complement and we know \CSIG\ is closed under union, then \CSIG\ would be
closed under complement (since $\overline{L_1} = \Sigma^* - L_1$, and
$\Sigma^*$ is CFL). But Lemma~9.4 shows that \CSIG\ is not closed under
complement. Hence \CSIG\ is \emph{not} closed under relative complement. \qed

%%------------------------------------------------------------
\item \textbf{$\CSIG$ is not closed under intersection or complement.}

\begin{enumerate}
\item Let $L_1 = \{\pmb{a}^i\pmb{b}^i\pmb{c}^j \mid i,j \in \NN\}$ and $L_2 = \{\pmb{a}^i\pmb{b}^j\pmb{c}^j \mid i,j \in \NN\}$. Both are context free. Their intersection:
\[
L_1 \cap L_2 = \{\pmb{a}^n\pmb{b}^n\pmb{c}^n \mid n \in \NN\},
\]
which is not context free (by the pumping lemma). So \CSIG\ is not closed under intersection.

If \CSIG\ were closed under complement, then $\overline{L_1}$ and $\overline{L_2}$ would be context free, and $\overline{\overline{L_1} \cup \overline{L_2}} = L_1 \cap L_2$ would be context free (by De Morgan and closure under union and complement). But $L_1 \cap L_2$ is not CFL. Contradiction. So \CSIG\ is not closed under complement.

\item For $|\Sigma| > 1$, choose two distinct symbols of $\Sigma$ and call them
$\pmb{a}$ and $\pmb{b}$.  Consider the languages
\[
L_1=\{\pmb{a}^i\pmb{b}^j\pmb{a}^j\mid i,j\in\NN\},\qquad
L_2=\{\pmb{a}^i\pmb{b}^i\pmb{a}^j\mid i,j\in\NN\}.
\]
Both are context free: $L_1$ is generated by
$S\to AB$, $A\to\pmb{a}A\mid\lambda$, $B\to\pmb{b}B\pmb{a}\mid\lambda$, and
$L_2$ is generated by
$S\to CD$, $C\to\pmb{a}C\pmb{b}\mid\lambda$, $D\to\pmb{a}D\mid\lambda$.

But
\[
L_1\cap L_2=\{\pmb{a}^n\pmb{b}^n\pmb{a}^n\mid n\in\NN\},
\]
which is not context free by the pumping lemma.  Therefore
$\mathfrak{C}_{\Sigma}$ is not closed under intersection whenever $|\Sigma|>1$.

If $\mathfrak{C}_{\Sigma}$ were closed under complement, then by De Morgan's law
it would also be closed under intersection, since it is already closed under
union.  This contradiction shows that $\mathfrak{C}_{\Sigma}$ is not closed
under complement either. \qed
\end{enumerate}

%%------------------------------------------------------------
\item \textbf{Automaton-based method for Theorem 9.3 (epsilon elimination).}

Consider a pushdown automaton (PDA) with states labeled by subsets $T \subseteq \Omega$ (the power set of nonterminals). Start state: $\{S\}$. Transitions: from state $T$, for each nonterminal $A \in T$ and production $A \to \gamma$, compute the set of nonterminals derivable in one step from $\gamma$, and union these into the new state.

This is similar to the subset construction for NFAs, adapted for CFGs. The set $T$ represents all nonterminals reachable in derivations. The nullable nonterminals ($N_\lambda$) are those states from which $\lambda$ is reachable.

%%------------------------------------------------------------
\item \textbf{Grammars in GNF with at most two symbols on the right.}

\begin{enumerate}
\item This is not a normal form for all CFLs: The language $\{\pmb{a}^n \mid n \geq 1\}$ can be generated (e.g., $S \to \pmb{a}S \mid \pmb{a}$), but more complex CFLs like $\{ww^r\}$ require productions of the form $S \to \pmb{a}S\pmb{a}$ in GNF which has 3 symbols on the right (more than 2). Converting to this restricted GNF may require more nonterminals and fail to represent the language if we strictly limit to 2 symbols. In fact $\pmb{a}Xa$ (where $a$ is terminal and $X$ is nonterminal) has 3 symbols; with at most 2 on the right in GNF (first terminal, at most one nonterminal), we can only write $\pmb{a}X$ (dropping the trailing $\pmb{a}$). So $\{ww^r \mid w \in \{\pmb{a},\pmb{b}\}^*\}$ cannot be generated.

\item Languages generated by this restricted form: GNF with at most two RHS symbols means productions are $A \to \pmb{a}$ or $A \to \pmb{a}B$. These are exactly the \emph{right-linear grammars}, which generate precisely the \emph{regular languages} (\DSIG\ = \GSIG\ = \RSIG). \qed
\end{enumerate}

%%------------------------------------------------------------
\item \textbf{Prove $L^*$ must be regular for \emph{any} collection $L$ over $\Sigma$.}

The statement is \emph{false} as written.

Take
\[
L=\{\pmb{a}^n\pmb{b}^n\mid n\geq 1\}.
\]
Then
\[
L^*=\{x_1x_2\cdots x_k\mid k\geq 0,\ x_i\in L\}.
\]
Assume for contradiction that $L^*$ is regular.  The language
$\pmb{a}^*\pmb{b}^*$ is regular, so
\[
L^*\cap \pmb{a}^*\pmb{b}^*
\]
would also be regular.

But
\[
L^*\cap \pmb{a}^*\pmb{b}^*=L.
\]
Indeed, any concatenation of two or more nonempty words from $L$ has at least
two alternations between $\pmb{a}$'s and $\pmb{b}$'s, so it cannot lie in
$\pmb{a}^*\pmb{b}^*$.  Thus the only words of $L^*$ that belong to
$\pmb{a}^*\pmb{b}^*$ are the single-block words $\pmb{a}^n\pmb{b}^n$.

Therefore $L$ would be regular, contrary to the standard result that
$\{\pmb{a}^n\pmb{b}^n\mid n\geq 1\}$ is not regular.  So $L^*$ need not be
regular.

%%------------------------------------------------------------
\item \textbf{If $|\Sigma|=1$, is $\CSIG$ closed under complementation?}

\textbf{Yes}.  If $|\Sigma|=1$, say $\Sigma=\{\pmb{a}\}$, and
$L\in\CSIG$, then Theorem~9.15 implies that $L$ is regular.  The regular
languages are closed under complementation, so $\overline{L}$ is also regular.
Every regular language is context free, hence $\overline{L}\in\CSIG$.

Therefore \CSIG\ is closed under complementation when $|\Sigma|=1$. \qed

%%------------------------------------------------------------
\item \textbf{$\{\pmb{a}^n\pmb{b}^n\pmb{c}^m \mid n,m \in \NN\}$ is context free.}

\textbf{Grammar}: $S \to AB$, $A \to \pmb{a}A\pmb{b} \mid \lambda$, $B \to \pmb{c}B \mid \lambda$.

$A$ generates $\{\pmb{a}^n\pmb{b}^n\}$ and $B$ generates $\pmb{c}^*$. So $L(G) = \{\pmb{a}^n\pmb{b}^n\pmb{c}^m\}$. \qed

%%------------------------------------------------------------
\item \textbf{Use Ogden's lemma to prove $\{\pmb{a}^k\pmb{b}^n\pmb{c}^m \mid k\neq n \wedge n\neq m\}$ is not context free.}

Assume $L = \{\pmb{a}^k\pmb{b}^n\pmb{c}^m \mid k\neq n, n\neq m\}$ is context free with Ogden constant $N$.

Choose $p > N$ and let $s = \pmb{a}^{p!+p}\pmb{b}^p\pmb{c}^{p!+p} \in L$ (since $p!+p \neq p$ and $p \neq p!+p$). Mark all $p$ occurrences of $\pmb{b}$ in $s$.

By Ogden's lemma: $s = uvwxy$ where $vwx$ contains at most $N \leq p$ marked positions, $vx$ contains at least 1 marked position, and all pumped strings $uv^i wx^i y \in L$.

Let $m = |vx|_{\pmb{b}} \geq 1$ (the number of $\pmb{b}$s in $vx$; at least 1 since $vx$ contains a marked symbol). Choose $i - 1 = p!/m$, which is a positive integer since $m \leq N < p$ implies $m \leq p-1$ and hence $m \mid p!$. The pumped string has $\pmb{b}$-count equal to $p + (i-1)m = p + p! = p!+p$.

Since $|vwx| \leq N < p$ and the $\pmb{b}$-block has exactly $p$ symbols, the window $vwx$ cannot reach both the $\pmb{a}$-block on the left and the $\pmb{c}$-block on the right simultaneously. Let $s_a = |vx|_{\pmb{a}}$ and $s_c = |vx|_{\pmb{c}}$. Then either $s_a = 0$ (window does not touch the $\pmb{a}$-block) or $s_c = 0$ (window does not touch the $\pmb{c}$-block).

\begin{itemize}
\item If $s_a = 0$: the $\pmb{a}$-count remains $p!+p = $ $\pmb{b}$-count, so $k = n = p!+p$, violating $k \neq n$. Contradiction.
\item If $s_c = 0$: the $\pmb{c}$-count remains $p!+p = $ $\pmb{b}$-count, so $n = m = p!+p$, violating $n \neq m$. Contradiction.
\end{itemize}

In every case the pumped string lies outside $L$, contradicting Ogden's lemma. \qed

\end{enumerate}
